\section{Background}
\label{sec:background}

\paragraph{DuckDB and its optimizer.} DuckDB~\cite{raasveldt2019duckdb} is an embedded analytical database written in
C++. A query passes through a parser, a binder, an optimizer and a vectorized executor. In DuckDB~2.0.0 the optimizer applies
a fixed sequence of 46 passes (for example filter pushdown, join ordering, removal of unused columns, top-N window
elimination). One pass, the \emph{expression rewriter}, applies small \emph{rules} (29 source files in
\texttt{src/optimizer/rule}) to every expression, for example
moving constants across a comparison or simplifying a \texttt{date\_trunc}. The setting
\texttt{disabled\_optimizers} switches off a named pass; a rule can only be switched off together with its whole pass.
We call this a \emph{switch}. DuckDB~2.0.0 was released during our project; we tested that release and checked every
candidate on the newest development build (\emph{main}).

\paragraph{Logic bugs and oracles.} A \emph{logic bug} makes the database return a wrong result without an error. To
detect one, a fuzzer needs an \emph{oracle}: a second computation whose result must be equal. Differential testing
compares two systems or two configurations of one system. Ternary logic partitioning (TLP)~\cite{rigger2020tlp} splits a
query by a predicate into three parts that must add up to the original. NoREC~\cite{rigger2020norec} compares an
optimized query with a form the optimizer cannot rewrite. Our main oracle is differential: the same query with a switch on
and off. An \emph{internal error} is DuckDB's own assertion that an impossible state was reached; a \emph{spurious error}
is an error raised for a valid query that should succeed.

\paragraph{Reach.} A test for a rule is useful only if the rule fires on it. We say a query \emph{reaches} a target
when its plan changes when the target's switch is turned off (\emph{expression-level reach}), and, more strictly, when a
call that only the target rule removes disappears from the plan (\emph{rule-level reach}). The fraction of valid queries
that reach the target is the \emph{reach rate}.

\paragraph{Terms used later.} A \emph{generator} is a short Python program that returns test cases (a setup script and
queries); an LLM writes it, our harness runs it. A \emph{hit} is a raw difference between configurations. A
\emph{candidate} is a hit that survived automatic triage. An \emph{issue} is a candidate a person checked and filed upstream.
